\section{Related research and contribution}
Research on lyrics has already documented long-term changes in popular music. Brand, Acerbi and Mesoudi studied emotional expression over five decades of song lyrics, showing how popular songs can be used to examine cultural transmission and emotional change~\cite{brand}. Their emotional measures do not by themselves distinguish romance, sexual content, wealth, weapons or the stance taken toward violence. Our study makes those distinctions explicit and compares their circulation across markets.

Varnum and colleagues analyzed 14,661 songs in US charts between 1958 and 2016. They linked simpler lyrics to a larger supply of new songs and to chart success~\cite{varnum}. Parada-Cabaleiro and colleagues examined changes in lyrical complexity, repetition and emotional language across genres~\cite{parada}. Together, these studies show why an aggregate change in lyrics needs to be examined alongside the supply and type of songs entering the charts. We address that issue through genre and release-group decompositions, while recognizing that broad artist-derived tags and incomplete dates cannot fully explain repertoire selection.

International listening also varies with cultural and social context. Liu, Hu and Schedl studied cross-country music preferences in relation to culture, socioeconomic conditions and friendship~\cite{liu}. Schedl and colleagues identified country archetypes from listening patterns and used them in music recommendation~\cite{schedl}. Those studies motivate regional comparisons, but listening-pattern clusters and clusters of lyrical themes measure different things. A country can share themes with another without sharing the same songs or listeners.

The added contribution is a common 26-dimension lyrical taxonomy linked to recording IDs, circulation markets, years, chart days, recorded streams and dated social indicators. It separates romantic affection, sexual explicitness, gender representation and material or violent themes instead of combining them into one score. It also compares counting songs with counting chart exposure, and tests how much temporal change is associated with genre and release-group composition. This is a descriptive extension of earlier work, not evidence that chart lyrics directly measure national values. Its stronger multilingual claims require the independent human validation described below.
